%=====================================================================
\chapter{Introduction}
%=====================================================================

The price of a traded security is the outcome of a continuous negotiation between
participants who do not observe the same information at the same time. Some of that
information is structured and numerically clean: open, high, low, close and volume
records, and the technical indicators derived from them. A larger part is unstructured
and arrives irregularly, in news articles, corporate disclosures, regulatory
announcements and social media commentary. A third part changes slowly and shapes the
environment rather than the individual security, in the form of policy rates, inflation,
industrial output and currency movements. Any forecasting system that reads only one of
these streams is, by construction, reasoning from an incomplete description of the
market it is trying to anticipate~\cite{snasel2024,long2024}.

Two developments have made the integration of these streams practical rather than
aspirational. The first is the maturation of multi-source data fusion, which supplies a
principled account of how heterogeneous streams may be combined so that their
complementary content is preserved instead of being averaged away~\cite{snasel2024,
zhang2022fusion}. The second is the arrival of transformer-based large language models,
which read financial text with a level of contextual sensitivity that lexicon and
bag-of-words methods never achieved, and which can therefore convert narrative
information into a numerically usable signal~\cite{vaswani2017,araci2019,li2024llmsurvey}.
Together they make it reasonable to ask a question that was previously out of reach:
whether a single system can absorb quantitative, textual and macroeconomic evidence and
convert it into a decision that an investor could actually execute.

That question is sharper than it appears, because three difficulties sit between a
fused signal and a defensible decision. The first is that the reliability of a source is
not constant. Technical indicators are informative in quiet trending conditions and
degrade when the market is disturbed; news and volatility signals behave in the opposite
manner. A fusion rule with fixed weights therefore misweights its inputs precisely when
accurate weighting matters most~\cite{farimani2024,tai2022}. The second is that a
prediction produced by a deep architecture is not, by itself, an account of why the
prediction was made. In a domain where capital is committed on the basis of the output,
opacity is a practical obstacle and not merely a philosophical one~\cite{koa2024,
lundberg2017}. The third is that predictive accuracy and economic usefulness are
different properties. A forecast becomes an investment only after it has been reconciled
with a covariance structure, a risk budget, position limits and the cost of trading, and
that translation can destroy the value of a forecast as easily as preserve
it~\cite{demiguel2009,gu2020}.

This research is directed at those three difficulties simultaneously. Its title,
\emph{Integrating Multi-Source Data with Sentiment Analysis and Language Models to
Enhance Stock Market Decision Making}, states the ambition, and the emphasis falls on
the final two words. The work is not an attempt to add another architecture to a crowded
field of predictors. It is an attempt to build, instrument and honestly evaluate the
complete path from heterogeneous raw information to an executable allocation, and to
determine which stages of that path actually contribute to the outcome.

Figure~\ref{fig:context} places the four research objectives on that path. Objective~1
governs the integration of sources. Objective~2 governs the extraction of signal from
language and its combination with conventional indicators. Objective~3 governs the
transparency of the reasoning and the behaviour of sentiment as a noisy, time-lagged
input. Objective~4 governs the conversion of a signal into a position and the validation
of that conversion across markets. The objectives are sequential in the sense that each
consumes the output of the one before it, but they are not merely stages of a pipeline:
each raises a distinct research question and each is answered by distinct experimental
evidence.

\begin{figure}[!htb]
\centering
\scalebox{0.86}{\begin{tikzpicture}[node distance=3.2mm]
\node[blk,minimum width=23mm] (m) {Market data\\OHLCV, indicators};
\node[blk,right=of m,minimum width=23mm] (n) {News and social\\media text};
\node[blk,right=of n,minimum width=23mm] (mac) {Macroeconomic\\series};
\node[blk,right=of mac,minimum width=23mm] (alt) {Alternative data\\volatility index};

\node[blkg,below=9mm of n,minimum width=52mm,xshift=9mm] (fuse)
  {\textbf{Objective 1} --- multi-source integration\\
   time-varying source weighting};
\node[blka,below=7mm of fuse,minimum width=52mm] (llm)
  {\textbf{Objective 2} --- LLM sentiment extraction\\
   hybrid fusion with financial indicators};
\node[blka,below=7mm of llm,minimum width=52mm] (xai)
  {\textbf{Objective 3} --- transparent reasoning\\
   sentiment noise and temporal dynamics};
\node[blkg,below=7mm of xai,minimum width=52mm] (alloc)
  {\textbf{Objective 4} --- validated decision making\\
   allocation, risk budget, cross-market test};
\node[blkd,below=7mm of alloc,minimum width=52mm] (dec)
  {Executable position with a stated risk profile};

\draw[ar] (m.south) -- ++(0,-3mm) -| (fuse.north west);
\draw[ar] (n.south) -- ++(0,-3mm) -| ([xshift=-14mm]fuse.north);
\draw[ar] (mac.south) -- ++(0,-3mm) -| ([xshift=14mm]fuse.north);
\draw[ar] (alt.south) -- ++(0,-3mm) -| (fuse.north east);
\draw[ar] (fuse) -- (llm);
\draw[ar] (llm) -- (xai);
\draw[ar] (xai) -- (alloc);
\draw[ar] (alloc) -- (dec);
\node[lbl,right=5mm of fuse,align=left] {what to fuse};
\node[lbl,right=5mm of llm,align=left] {how to read text};
\node[lbl,right=5mm of xai,align=left] {why to trust it};
\node[lbl,right=5mm of alloc,align=left] {how much to hold};
\end{tikzpicture}}
\caption[Research context and the position of the four objectives]{Research context: the path from heterogeneous market information to an
executable decision, and the position of the four research objectives on that path.}
\label{fig:context}
\end{figure}

The Indian equity market is the primary empirical setting. The choice is deliberate
rather than convenient. The National Stock Exchange combines high retail participation,
strong sensitivity to domestic policy events such as the Union Budget and Reserve Bank of
India monetary policy announcements, pronounced sectoral concentration in financial
services, and simultaneous exposure to global factors including crude oil prices and
United States monetary policy~\cite{mu2023investor,bacco2024}. These features make
it an unusually demanding test of a multi-source system, because the sources genuinely
disagree and the disagreement is informative. Cross-market validation on United States
multi-asset and mega-capitalisation universes is carried out separately under
Objective~4 in order to establish where the framework transfers and where it does not.

The remainder of this report is organised as follows. Chapter~2 reviews the literature
on multi-source fusion, sentiment analysis with language models, explainable and causal
financial modelling, and decision-level evaluation, and states the research gaps that
follow from it. Chapter~3 formalises the problem statement and the four research
objectives. Chapter~4 states the research contributions and maps the candidate's
published work onto those objectives. Chapter~5 presents the methodology, the
implementation and the experimental evidence for each objective, and then shows how the
individual contributions compose into a single system. Chapter~6 discusses the
consolidated results and the significance of the work. Chapter~7 concludes and states
the future scope, and Chapter~8 gives the intended organisation of the thesis.

%=====================================================================
\chapter{Literature Survey}
%=====================================================================

This chapter reviews the research directly relevant to the thesis. It is organised along
the same four axes as the objectives, so that the gap identified at the end of each
section is the gap that the corresponding objective addresses. Emphasis is placed on
what each body of work establishes and, more importantly, on what it leaves unresolved.

\section{Multi-Source Data Fusion in Financial Prediction}

The premise that heterogeneous sources outperform any single source is now well
supported. Sn\'a\v{s}el et al.~\cite{snasel2024} formalised a generalised multi-source
fusion framework for stock selection and demonstrated that complementary streams
mitigate the individual weaknesses of each, establishing the conceptual basis for
combining numerical market data with unstructured text. Long et al.~\cite{long2024}
reached the same conclusion from a multi-view heterogeneous data perspective, showing
that different views of the same security carry non-overlapping information. Mu et
al.~\cite{mu2023investor} combined investor sentiment with optimised deep learning and
reported consistent gains over price-only baselines, and Shaban et al.~\cite{shaban2024}
confirmed the pattern for trend forecasting. Parallel work on single-source architectures
establishes the ceiling that fusion must beat: Song~\cite{song2024} enriched the feature
space for trend-reversal prediction and Javaid et al.~\cite{javaid2024} combined a peephole
recurrent network with a time-attention layer, both improving on conventional deep
baselines while remaining confined to price-derived inputs.

The methodological question is not whether to fuse but how. Three strategies dominate.
Early fusion concatenates features from all sources before modelling, which is simple but
forces one architecture to learn cross-domain interactions and is prone to modality
dominance. Late fusion trains source-specific models and combines their predictions,
preserving modularity but usually fixing the combination weights in advance. Hybrid or
attention-based fusion learns modality importance inside the
network~\cite{choi2023,haryono2023}, a genuine advance that nonetheless produces weights
which are point estimates over an entire training set rather than statements about the
reliability of a source at a particular moment. Moodi et al.~\cite{moodi2024} concatenate technical indicators with sentiment
scores inside a hybrid convolutional and recurrent model, Liu and
Wang~\cite{liu2019attention} weight features by a numerical attention mechanism, and
Zong et al.~\cite{zong2025} stabilise multimodal fusion with gated cross-attention.
Farimani et al.~\cite{farimani2024} proposed adaptive multimodal weighting and
Tai et al.~\cite{tai2022} modelled the joint distribution of sentiment and financial
signals with normalising flows; both move toward adaptivity without an explicit,
probabilistic measure of source confidence. Where fusion output is carried through to a
portfolio at all, the conversion is typically heuristic: Sami et al.~\cite{sami2025}
construct holdings by statistical averaging of forecasts, without a covariance structure
or an explicit risk objective.

\section{Sentiment Analysis and Language Models in Finance}

Financial text analysis has passed through three generations. Dictionary methods scored
text against predefined polarity lexicons; Loughran and McDonald~\cite{loughran2011}
showed that general-purpose lexicons misclassify financial language badly enough to
require domain-specific alternatives, and VADER~\cite{hutto2014} remains a widely used
reference implementation. Supervised machine learning improved accuracy but required large
annotated corpora and generalised poorly across market conditions~\cite{rouf2021}.
Transformer architectures~\cite{vaswani2017} displaced both: FinBERT~\cite{araci2019}
adapted bidirectional encoders to financial language, and domain models including
BloombergGPT~\cite{wu2023bloomberggpt} and the FinGPT
family~\cite{yang2023fingpt,wang2024fingptbench} extended this to generative reasoning
over financial documents. Surveys~\cite{li2024llmsurvey,todd2024,bahoo2024} converge on
the finding that contextual language models extract market-relevant meaning that lexicon
methods miss.

Applications have followed quickly. Fatouros et al.~\cite{fatouros2024} showed that
conversational language models outperform established financial sentiment baselines,
Wilksch and Abramova~\cite{wilksch2024} built a dedicated model for sentiment in financial
tweets, and Khan and Umer~\cite{khan2024chatgpt} surveyed the applications and the
practical obstacles. Haryono et al.~\cite{haryono2023} coupled a transformer news encoder
with a gated recurrent unit for technical data and later extended the approach to
generative aspect-based sentiment~\cite{haryono2024}, and Ho and Huang~\cite{ho2021}
combined sentiment with candlestick chart representations. Jung and
Jang~\cite{jung2024} improved financial sentiment quality by masking numerically
significant tokens so that quantitative context is preserved, and \'On\'oz\'o et
al.~\cite{onozo2024} applied language models to news analysis and macroeconomic
nowcasting. Applied studies span markets and asset classes: Vargas et
al.~\cite{vargas2022} fused sentiment with recurrent price models on Brazilian equities,
Wang and Chen~\cite{wang2024spcm} built a sentiment-based recommendation system, and
Shang~\cite{shang2025} combined high-dimensional indicators with social media dynamics.
Elahi and Taghvaei~\cite{elahi2024} combined financial data with news articles using
large language models directly, and Mishra et al.~\cite{mishra2025} surveyed the
resulting body of work. What this literature establishes is that language models improve
sentiment measurement. What it establishes far less clearly is how much of that
improvement survives into forecast quality once the sentiment score is fused with price
history, aligned to trading sessions and evaluated without look-ahead.

\section{Explainability, Causality and Uncertainty}

Opacity is a recognised obstacle to the adoption of learned models in finance. Koa et
al.~\cite{koa2024} used self-reflective large language models to generate natural-language
explanations alongside stock predictions, and the general-purpose attribution methods of
Lundberg and Lee~\cite{lundberg2017} and Ribeiro et al.~\cite{ribeiro2016} are now
standard instruments for post-hoc interpretation. Post-hoc attribution, however, explains
a model; it does not establish that the relationships the model has learned are anything
more than correlations that happen to hold in the training window.

Causal formulations address this directly. Xu et al.~\cite{xu2026causal} showed that
causal feature selection is more robust out of sample than correlation-based feature
engineering, but their causal graph is static and cannot represent the structural breaks
produced by policy announcements. Liu et al.~\cite{liu2026dynamic} introduced a
dynamic-causal lens that encodes stock interactions as symbolic movement patterns and
adapts to changing conditions, at a computational cost that scales quadratically in the
number of instruments. Yu and Liu~\cite{yu2025causalseqgnn} coupled factor investing with
causal graph neural networks using Fama-French factors calibrated to United States
equities. Luo et al.~\cite{luo2025magicnet} added memory to a causal interaction network,
and Xing et al.~\cite{xing2024vague} modelled the vagueness of inter-stock relations.
Alongside this, multi-granularity architectures~\cite{zhu2025multigran,wang2025mtmd}
demonstrated that financial series carry structure at several time scales at once, with
Wang et al.~\cite{wang2025mtmd} adding a debiasing layer that assumes stationary noise
properties.

Relational formulations extend the same idea to the cross-section. Patel et
al.~\cite{patel2024hybrid} treat securities as nodes in a learned relational graph, Liu
et al.~\cite{liu2024heterogeneous} model mutual interplay through dual-dynamic attention,
and Gao et al.~\cite{gao2025relational} select stocks through probabilistic state-space
learning. The statistical instrument underlying directed inference in this setting
remains the notion of predictive precedence formalised by
Granger~\cite{granger1969} and the constraint-based discovery procedures of Spirtes et
al.~\cite{spirtes2000}, while the learning components themselves rest on long
short-term memory~\cite{hochreiter1997} and gradient-boosted tree
ensembles~\cite{chen2016xgboost}.

A related and less examined issue is that a signal's reliability is itself uncertain.
Most fusion architectures pass point predictions forward without any statement of
confidence, so a source that has recently become noisy continues to be weighted as though
it had not. Explicit, time-varying uncertainty quantification is the natural remedy and
is largely absent from the reviewed fusion literature.

\section{From Prediction to Decision: Allocation and Validation}

Portfolio construction has its own long literature. Markowitz~\cite{markowitz1952}
established the mean-variance formulation, whose extreme sensitivity to estimated inputs
motivated both shrinkage estimators~\cite{ledoit2004} and the demonstration by DeMiguel et
al.~\cite{demiguel2009} that naive equal weighting can outperform optimisation that
ignores estimation error. Regime-switching models following Hamilton~\cite{hamilton1989}
showed that financial series are not homogeneous through time; Ang and
Bekaert~\cite{ang2002} demonstrated that accounting for regime shifts materially changes
optimal allocations, and Kritzman et al.~\cite{kritzman2012} made the practical case for
regime-aware dynamic strategies. Volatility-managed portfolios~\cite{moreira2017} scale
total risk by recent volatility but say nothing about composition. Kumar and
Yadav~\cite{kumar2024rdmdea} proposed credibilistic portfolio selection under
RDM-DEA, handling positive and negative returns without assuming normality, but without
integration to a dynamic signal. Risk-based construction, of which the hierarchical
approach of L\'opez de Prado~\cite{lopezdeprado2016} is representative, forgoes signal
information by design. Gu et al.~\cite{gu2025ragic} generated risk-aware prediction
intervals, and Zhu et al.~\cite{zhu2025spin} built sparse portfolios driven by irregular
news streams using English-language sources only.

A second and older family of sentiment measures reads sentiment off the market rather
than off text. Baker and Wurgler~\cite{bakerwurgler2007} construct investor sentiment
from trading activity and dispersion, on the argument that participation reveals what
surveys merely report. Singsiri and Yokrattanasak~\cite{singsiri2026} build allocation
signals from money-flow indicators on exactly this reasoning, and Naidoo et
al.~\cite{naidoo2025} model investor sentiment and market volatility as a joint object.
This family matters here because instrument-level news corpora do not exist for broad
index instruments, which constrains what can be measured at the allocation layer.

Evaluation practice is the final concern. Gu et al.~\cite{gu2020} provide the benchmark
machine-learning study in asset pricing and note that measured predictability is small.
Harvey and Liu~\cite{harvey2015} showed that conventional significance thresholds are
inadequate once the number of strategies examined is taken into account. Robust inference
under serial dependence requires corrections such as those of Newey and
West~\cite{newey1987} and resampling schemes such as the stationary
bootstrap~\cite{politis1994}. The practical consequence is that a reported gain is
credible only when the specification search behind it is reported alongside it.

Table~\ref{tab:lit} consolidates the representative work reviewed above, stating for each
the contribution that is relevant to this thesis and the limitation that this thesis must
work around.

\begin{table}[!htb]
\centering
\caption[Consolidated review of representative work]{Consolidated review of representative work, with the limitation relevant to
this research}
\label{tab:lit}
\tabfont
\begin{tabular}{@{}L{2.35cm}L{3.5cm}L{3.9cm}L{4.3cm}@{}}
\toprule
\textbf{Reference} & \textbf{Technique} & \textbf{Contribution relevant here} &
\textbf{Limitation to be addressed} \\
\midrule
Sn\'a\v{s}el et al.~\cite{snasel2024} & Generalised multi-source fusion for stock
selection & Formal basis for combining numerical and textual streams & Fusion weights are
not conditioned on the prevailing market state \\
Long et al.~\cite{long2024} & Multi-view heterogeneous data model & Views carry
non-overlapping information & Developed and tested outside emerging-market conditions \\
Haryono et al.~\cite{haryono2023}; Mu et al.~\cite{mu2023investor} & Sentiment fused with
technical indicators & Gated and optimised fusion of text and price streams & Gate learned
globally; sentiment quality assumed rather than validated \\
Farimani et al.~\cite{farimani2024} & Adaptive multimodal weighting & Moves toward
time-varying source importance & No explicit probabilistic confidence per source \\
Araci~\cite{araci2019}; Wu et al.~\cite{wu2023bloomberggpt} & Domain-adapted language
models & Contextual sentiment extraction from financial text & Gain in sentiment quality
is not traced through to decision quality \\
Koa et al.~\cite{koa2024} & Self-reflective LLM explanations & Natural-language
justification of predictions & Explains the model, not the causal structure \\
Xu et al.~\cite{xu2026causal}; Liu et al.~\cite{liu2026dynamic} & Causal and
dynamic-causal graph modelling & Causal features are more robust out of sample &
Static graph cannot absorb policy shocks; dynamic variant scales quadratically \\
Wang et al.~\cite{wang2025mtmd} & Multi-scale temporal memory with debiasing & Multiple
horizons processed jointly & Debiasing assumes stationary noise \\
Kumar and Yadav~\cite{kumar2024rdmdea} & RDM-DEA credibilistic portfolio selection &
Optimisation without a normality assumption & Not coupled to a dynamic signal for
rebalancing \\
Moreira and Muir~\cite{moreira2017} & Volatility-managed portfolios & Total risk scaled
by recent volatility & Does not address portfolio composition \\
DeMiguel et al.~\cite{demiguel2009}; Harvey and Liu~\cite{harvey2015} & Naive
diversification; backtesting under multiple testing & Estimation error can defeat
optimisation; reported gains must account for specification search & Sets the benchmark
any framework must beat and requires every specification to be disclosed \\
\bottomrule
\end{tabular}
\end{table}

\section{Research Gaps}

The review supports five gaps, which are summarised in Figure~\ref{fig:taxonomy} and
stated below.

\begin{enumerate}[label=\textbf{G\arabic*.}]
\item \textbf{Static fusion of dynamically reliable sources.} Existing frameworks fuse
financial, textual and macroeconomic streams using fixed weights, globally learned
attention or simple averaging. None conditions the contribution of a source on a
measured, time-varying estimate of that source's own predictive reliability, so a stream
that has become noisy continues to influence the forecast at its historical weight.

\item \textbf{Incomplete coverage of the Indian information environment.} Most
multi-source frameworks are designed and validated on developed markets. Frameworks that
jointly incorporate financial indicators, news sentiment, social media signals and Indian
macroeconomic variables --- policy rate, inflation, industrial production, currency ---
under a single architecture and a strict temporal protocol are rare.

\item \textbf{Unverified translation from language-model sentiment to decision quality.}
Language models demonstrably improve sentiment measurement. Whether that improvement
survives temporal alignment, fusion with technical indicators and honest out-of-sample
evaluation at a single-name horizon has not been established, and reported gains are
seldom accompanied by the baseline against which they should be judged.

\item \textbf{Correlational rather than mechanistic transparency.} Post-hoc attribution
explains a fitted model. It does not distinguish a genuine driver from a spurious
association, and it does not describe how a sentiment shock propagates into a price
movement over time. A transparency mechanism grounded in directed causal structure and in
multi-scale temporal dynamics is required, together with a direct measurement of how much
noise the sentiment channel actually carries.

\item \textbf{Prediction reported in place of decision.} Statistical accuracy remains the
dominant success criterion, and economic value is reported inconsistently. Where
allocation is addressed it is often heuristic, and the contribution of individual
architectural layers is almost never isolated, so architectures accumulate components
whose value has never been tested. Cross-market transfer is asserted more often than it
is demonstrated.
\end{enumerate}

\begin{figure}[!htb]
\centering
\scalebox{0.86}{\begin{tikzpicture}[node distance=4mm]
\node[blkd,minimum width=32mm] (root) {Multi-source financial\\prediction and decision};

\node[blk,below left=5mm and 22mm of root,minimum width=25mm,minimum height=6mm] (a) {Fusion strategy};
\node[blk,below=5mm of root,minimum width=25mm,xshift=-2mm,minimum height=6mm] (b) {Signal extraction};
\node[blk,below right=5mm and 20mm of root,minimum width=25mm,minimum height=6mm] (c) {Decision layer};

\node[blk,below=4mm of a,minimum width=27mm,minimum height=6mm] (a1)
  {early / late / attention\\\textit{fixed or global weights}};
\node[blk,below=4mm of b,minimum width=27mm,minimum height=6mm] (b1)
  {lexicon $\to$ ML $\to$ LLM\\\textit{measurement improves}};
\node[blk,below=4mm of c,minimum width=27mm,minimum height=6mm] (c1)
  {accuracy reported\\\textit{allocation heuristic}};

\node[blka,below=4mm of a1,minimum width=27mm,minimum height=8mm] (a2)
  {\textbf{G1, G2}\\uncertainty-weighted,\\market-specific fusion};
\node[blka,below=4mm of b1,minimum width=27mm,minimum height=8mm] (b2)
  {\textbf{G3, G4}\\validated sentiment,\\causal transparency};
\node[blka,below=4mm of c1,minimum width=27mm,minimum height=8mm] (c2)
  {\textbf{G5}\\ablated allocation,\\cross-market test};

\draw[ar] (root) -- (a); \draw[ar] (root) -- (b); \draw[ar] (root) -- (c);
\draw[ar] (a) -- (a1); \draw[ar] (b) -- (b1); \draw[ar] (c) -- (c1);
\draw[arb] (a1) -- (a2); \draw[arb] (b1) -- (b2); \draw[arb] (c1) -- (c2);
\node[lbl,rotate=90,anchor=center] at ([xshift=-8mm]a1.west) {state of the art};
\node[lbl,rotate=90,anchor=center] at ([xshift=-8mm]a2.west) {this research};
\end{tikzpicture}}
\caption[Taxonomy of the reviewed field and the five research gaps]{Taxonomy of the reviewed field and the location of the five research gaps
addressed by this work.}
\label{fig:taxonomy}
\end{figure}

%=====================================================================
\chapter{Problem Statement and Research Objectives}
%=====================================================================

\section{Problem Statement}

To design, develop and validate an integrated decision-support framework that fuses
heterogeneous financial information --- market metrics, news and social media sentiment
extracted using large language models, and macroeconomic indicators --- into
interpretable and reliability-aware predictions, and that converts those predictions into
executable, risk-controlled stock market decisions whose performance is demonstrated
across market states and across markets. The framework must address the limitations of
existing approaches, namely static source weighting under non-stationary reliability,
opaque and correlational reasoning, unquantified sentiment noise, and the routine
substitution of statistical accuracy for economic value.

\section{Research Objectives}

\begin{enumerate}[label=\textbf{RO\arabic*.},leftmargin=2.6em]
\item To investigate and develop a comprehensive model that integrates diverse data
sources --- such as financial metrics, news articles, social media sentiment and
macroeconomic factors --- aiming to enhance the accuracy of stock market predictions by
providing a holistic analysis.

\item To design a hybrid framework that combines large language models, sentiment
analysis techniques and conventional financial indicators, focusing on optimising deep
learning architectures for real-time stock market predictions.

\item To establish and evaluate an explainable artificial intelligence framework capable
of delivering transparent and interpretable predictions in stock selection, analysing the
impact of sentiment data noise and exploring the temporal dynamics between sentiment
shifts and stock price movements.

\item To evaluate and validate the result and practical application of the proposed model
across different stock markets, exploring transfer learning techniques for cross-market
applications and validating the model against benchmark datasets and real-world market
data.
\end{enumerate}

Figure~\ref{fig:gapmap} traces each gap identified in Section~2.5 to the objective that
addresses it and to the contribution reported in Chapter~4. The mapping is
many-to-many by design: G1 and G2 are jointly answered by the fusion work of RO1, while
G5 requires evidence from both the decision layer of RO4 and the transparency instruments
of RO3.

\begin{figure}[!htb]
\centering
\scalebox{0.86}{\begin{tikzpicture}[node distance=3mm]
\node[blkd,minimum width=30mm,minimum height=6mm] (g1) {\textbf{G1} static fusion weights};
\node[blkd,below=3mm of g1,minimum width=30mm,minimum height=6mm] (g2) {\textbf{G2} Indian information environment};
\node[blkd,below=3mm of g2,minimum width=30mm,minimum height=6mm] (g3) {\textbf{G3} sentiment $\to$ decision unverified};
\node[blkd,below=3mm of g3,minimum width=30mm,minimum height=6mm] (g4) {\textbf{G4} correlational transparency};
\node[blkd,below=3mm of g4,minimum width=30mm,minimum height=6mm] (g5) {\textbf{G5} prediction reported as decision};

\node[blkg,right=26mm of g1,minimum width=27mm,minimum height=6mm] (o1) {\textbf{RO1} integrated multi-source model};
\node[blkg,below=5mm of o1,minimum width=27mm,minimum height=6mm] (o2) {\textbf{RO2} hybrid LLM + indicator framework};
\node[blkg,below=5mm of o2,minimum width=27mm,minimum height=6mm] (o3) {\textbf{RO3} explainable and causal framework};
\node[blkg,below=5mm of o3,minimum width=27mm,minimum height=6mm] (o4) {\textbf{RO4} validation and cross-market transfer};

\node[blka,right=24mm of o1,minimum width=27mm,minimum height=6mm] (c1) {\textbf{RC1} MSFN and Bayesian dynamic fusion};
\node[blka,below=5mm of c1,minimum width=27mm,minimum height=6mm] (c2) {\textbf{RC2} LLM sentiment--indicator pipeline};
\node[blka,below=5mm of c2,minimum width=27mm,minimum height=6mm] (c3) {\textbf{RC3} DCHF and sentiment diagnostics};
\node[blka,below=5mm of c3,minimum width=27mm,minimum height=6mm] (c4) {\textbf{RC4} decision-theoretic allocation};

\draw[ar] (g1) -- (o1); \draw[ar] (g2) -- (o1);
\draw[ar] (g3) -- (o2); \draw[ar] (g4) -- (o3);
\draw[ar] (g5) -- (o4); \draw[arb] (g5) -- (o3);
\draw[ar] (o1) -- (c1); \draw[ar] (o2) -- (c2);
\draw[ar] (o3) -- (c3); \draw[ar] (o4) -- (c4);
\node[lbl,above=1.5mm of g1] {Research gaps};
\node[lbl,above=1.5mm of o1] {Research objectives};
\node[lbl,above=1.5mm of c1] {Research contributions};
\end{tikzpicture}}
\caption[Alignment of research gaps, objectives and contributions]{Alignment of research gaps, research objectives and research contributions.
Dashed links denote partial rather than primary coverage.}
\label{fig:gapmap}
\end{figure}
